\section{Preregistered evaluation protocol}

The authoritative registration is in \texttt{preregistration/}. The full configuration is the v0.1.0 file under \texttt{registered-configs/}; the smoke configuration is separate and is not a substitute for the campaign. The primary comparison is memory quality under matched task, history, model, and resource conditions.

\subsection{Primary memory baselines}

\scriptsize
\begin{longtable}{@{}p{0.16\linewidth}p{0.74\linewidth}@{}}
\toprule
ID & Definition \\
\midrule
M0 & No persistent memory; current task context only. \\
M1 & Full-history memory. \\
M2 & Periodically summarized textual memory. \\
M3 & Vector-similarity retrieval over historical observations. \\
M4 & Structured key/value or episodic memory. \\
M5 & Knowledge-graph memory where practical. \\
M6 & Full CCM: bounded local models, frames, history, recall, reconciliation, and optional active evidence. \\
M7 & CCM plus an optional ontology substrate; exploratory if implemented. \\
\bottomrule
\end{longtable}
\normalsize

The deterministic v0.1.0 adapter implements transparent symbolic controls for M0--M6. It does not claim to be a learned summarizer, embedding model, or production knowledge graph. The full registered campaign should use explicit model/runtime implementations where resources permit.

\subsection{CCM ablations}

\begin{longtable}{@{}p{0.25\linewidth}p{0.65\linewidth}@{}}
\toprule
Ablation & Removed or changed mechanism \\
\midrule
\texttt{CCM-NoRF} & Explicit reference-frame constraints. \\
\texttt{CCM-NoModules} & Multiple bounded columns; one local module remains. \\
\texttt{CCM-NoTemporal} & Temporal recall and current-state constraints. \\
\texttt{CCM-NoReconciliation} & Reconciliation-driven structured commit. \\
\texttt{CCM-NoActive} & Active evidence acquisition. \\
\texttt{CCM-FlatLedger} & Frame-aware organization; flat structured-ledger comparison. \\
\bottomrule
\end{longtable}

\subsection{Memory task classes}

The deterministic suite prioritizes persistent factual recall, temporal state, relational recall, conflicting evidence, state replacement/supersession, provenance recall, cross-frame integration, distractor resistance, long-horizon continuity, active evidence, incomplete information, and pure reasoning control. Each task has a stable ID, hidden answer, candidate hypotheses, explicit observations or actions, and required evidence where applicable.

\subsection{Metrics}

Primary metrics are recall precision, recall recall, state reconstruction accuracy, temporal accuracy, relation accuracy, provenance accuracy, contradiction preservation, false-memory rate, stale-state rate, and relevant-context density. Secondary metrics include downstream task success, unsupported-memory rate, context tokens, retrieval/update latency, storage bytes, relation count, index size, candidates examined, communication, calibration, abstention, unresolved outcomes, and failure taxonomy. Smoke tables expose these fields but do not constitute campaign results.

\subsection{Model scale and composition}

Model scale is a secondary dimension. Where practical, the same memory architecture should be evaluated across small local, larger open, and hosted/frontier model classes, measuring relative improvement from each model's baseline memory condition rather than asking whether a small model catches a frontier model. Multi-agent composition is exploratory: independent CCM-equipped agents may exchange bounded memory-derived state without requiring a shared global store.

\subsection{Resource accounting}

The harness records model invocations, input/output tokens, wall-clock latency, communication, bytes stored, object and relation counts, index size, recall candidates examined, context tokens produced, retrieval latency, update cost, and unavailable hardware/energy/API fields. Pareto analysis therefore compares memory quality against storage, retrieval latency, context, and compute.

\subsection{Statistical decision rules}

The unit of analysis is a task instance, paired by task ID and seed. The registered configuration uses three independent seeds and ten instances per task class. The analysis reports raw differences, 95\% percentile bootstrap confidence intervals using 10,000 resamples, paired effect sizes, and Holm correction across confirmatory comparisons. A minimum practically meaningful difference is 0.05 on the relevant bounded metric. Failures and timeouts remain in denominators; there is no silent outlier removal or outcome-dependent stopping.

\subsection{Negative results and provenance}

All registered instances, unresolved states, contradictions, abstentions, timeouts, and unavailable resource measurements remain in the analysis. Run manifests record configuration hash, Git commit, model/task versions, seed, hardware, timestamp, and raw output. Private chain-of-thought is neither required nor stored. The protocol makes no claim until the full campaign replaces the smoke status with traceable results.
